% Part: normal-modal-logic
% Chapter: syntax-and-semantics
% Section: introduction

\documentclass[../../../include/open-logic-section]{subfiles}

\begin{document}

\olfileid{nml}{syn}{int}

\olsection{परिचय}

मोडल तर्कशास्त्र \emph{मोडल विधाने} आणि त्यांच्यातील तार्किक
निष्पन्नता-संबंध यांचा अभ्यास करते. मोडल विधानांची पुढील
उदाहरणे पाहा:
\begin{enumerate}
\item $2+2=4$ हे अनिवार्य आहे.
\item उद्या पाऊस पडणे अनिवार्यपणे संभाव्य आहे.
\item जर~$!A$ अनिवार्यपणे संभाव्य असेल, तर~$!A$ संभाव्य आहे.
\end{enumerate}
संभाव्यता आणि अनिवार्यता या एकमेव मोडलता नाहीत: इतर एकस्थानी
संयोजकही मोडलता मानले जातात. उदाहरणार्थ, ``$!A$ असे असायला हवे,''
``पुढे~$!A$ घडेल,'' ``Dana ला~$!A$ माहीत आहे,'' किंवा
``Dana चा~$!A$ वर विश्वास आहे.''

मोडल तर्कशास्त्र प्रथम अरिस्टॉटल यांच्या \emph{De
  Interpretatione} मध्ये दिसते. अनिवार्य असल्यास संभाव्य असते,
पण उलट नाही; संभाव्यता आणि अनिवार्यता एकमेकींच्या साहाय्याने
परिभाषित करता येतात; $!A \land !B$ संभाव्यतः सत्य असेल तर
$!A$ संभाव्यतः सत्य असते आणि $!B$ संभाव्यतः सत्य असते,
पण उलट नाही; तसेच $!A \to !B$ अनिवार्य असेल, तर $!A$
अनिवार्य असताना $!B$ ही अनिवार्य असते, या गोष्टी प्रथम
त्यांच्या लक्षात आल्या.

मोडल तर्कशास्त्राकडे पाहण्याचा पहिला आधुनिक दृष्टिकोन
C.~I. Lewis यांच्या कार्यातून पुढे आला. Lewis आणि Langford
यांच्या \emph{Symbolic Logic} (1932) मध्ये तो परिपूर्णतेला
पोहोचला. वास्तविक अभिव्यंजन वापरून अभिव्यंजनाचे प्रतिनिधित्व
करण्याबाबत Lewis आणि Langford समाधानी नव्हते: $!A \lif !B$
हे ``$!A$ पासून $!B$ निष्पन्न होते'' यासाठी अपुरे पर्याय आहे.
त्याऐवजी ``अनिवार्यतः, जर $!A$ तर $!B$'' असे अभिव्यंजनाचे
लक्षण मांडून त्यांनी ते $!A \strictif !B$ असे दर्शवले.
या विविध गुणधर्मांची मांडणी करताना Lewis यांनी पाच भिन्न
मोडल प्रणाली ओळखल्या: \Log{S1}, \ldots, \Log{S4}, \Log{S5};
त्यांतील शेवटच्या दोन अजूनही वापरात आहेत.

Lewis आणि Langford यांचा दृष्टिकोन पूर्णतः \emph{विन्यासमीमांसक}
होता: त्यांनी योग्य वाटणारी स्वयंसिद्धके व नियम निवडले आणि
त्यांच्या साहाय्याने काय सिद्ध करता येते याचा अभ्यास केला.
चिन्हार्थमीमांसक दृष्टिकोन बराच काळ उपलब्ध नव्हता.
नंतर Rudolf Carnap यांनी \emph{Meaning and Necessity} (1947)
मध्ये \emph{अवस्था-वर्णन} ही संकल्पना वापरून पहिला प्रयत्न केला.
अवस्था-वर्णन म्हणजे त्या वर्णनात ``सत्य'' असलेल्या आण्विक
वाक्यांचा संग्रह. कोणत्याही वाक्य~$!A$ साठी सत्यतेची व्याख्या
विस्तारित केल्यानंतर, $!A$ सर्व अवस्था-वर्णनांत सत्य असेल
तर ते \emph{अनिवार्यतः सत्य} आहे, अशी Carnap यांची व्याख्या आहे.
मात्र त्यांचा दृष्टिकोन \emph{पुनरावृत्त} मोडलता हाताळू शकत
नव्हता: ``संभाव्यतः अनिवार्यतः \ldots संभाव्यतः $!A$''
या रूपाची वाक्ये नेहमी सर्वांत आतील मोडलतेपुरती सीमित होत.

मोडल चिन्हार्थमीमांसेतील मोठी प्रगती Saul Kripke यांच्या
``A Completeness Theorem in Modal Logic'' (JSL 1959) या
लेखामुळे झाली. एखादे विधान ``सर्व संभाव्य जगांत'' सत्य
असेल तर ते अनिवार्यतः सत्य असते, या लाइब्निझ यांच्या कल्पनेवर
Kripke यांनी आपले कार्य आधारले. परंतु त्या साध्या अर्थलक्षणातही
Carnap यांच्या दृष्टिकोनासारखीच अडचण आहे: जग~$w$ येथे एखादे
विधान अनिवार्यतः सत्य आहे की नाही, हे $w$ च्या निवडीवर
(किंवा अवस्था-वर्णन~$s$ वर) मुळीच अवलंबून नसते.
म्हणून Kripke यांनी जगे \emph{प्राप्यता संबंध}~$R$ ने
संबंधित आहेत असे मानले; आणि ``अनिवार्यतः $!A$'' या रूपाचे
विधान जग~$w$ येथे तेव्हा आणि केवळ तेव्हाच सत्य असते,
जेव्हा $w$ पासून \emph{प्राप्य असलेल्या} प्रत्येक जग~$w'$
येथे $!A$ सत्य असते. या दृष्टिकोनाचे कोणतेही रूप वापरणाऱ्या
चिन्हार्थमीमांसेला क्रिप्के चिन्हार्थमीमांसा म्हणतात;
तिच्यामुळे मोडल तर्कशास्त्रांचा, अनेकवचनात, झपाट्याने विकास
होऊ शकला.

क्रिप्के चिन्हार्थमीमांसेनुसार अर्थ लावल्यास, मोडल तर्कशास्त्र
\emph{संबंधी रचना} ``आतून'' कशा दिसतात हे दाखवते. संबंधी
रचना म्हणजे द्विस्थानी संबंध असलेला संच; उदाहरणार्थ,
वर्गातील विद्यार्थ्यांचा सामाजिक सुरक्षा क्रमांकानुसार
लावलेला संच ही संबंधी रचना आहे. पण संबंधी रचना अनेक
क्षेत्रांत आढळतात. जगाच्या अवस्थांमधील सापेक्ष संभाव्यतेखेरीज,
एखाद्या व्यक्तीच्या ज्ञानाच्या अवस्था ज्ञानविषयक संभाव्यतेने
संबंधित असू शकतात किंवा गतिशील प्रणालीच्या अवस्थांमध्ये
स्थित्यंतरे असू शकतात. या सर्वांचे प्रतिरूपण मोडल तर्कशास्त्राने
करता येते: पहिल्या प्रकारातून नेहमीचे सत्यविषयक मोडल तर्कशास्त्र,
तर इतर प्रकारांतून ज्ञानविषयक तर्कशास्त्र, गतिशील तर्कशास्त्र
इत्यादी मिळतात.

मोडल तर्कशास्त्रज्ञ ``अनुरूपता सिद्धांत'' म्हणतात त्या एका
विशिष्ट दृष्टीकोनावर आपण लक्ष केंद्रित करतो. Kripke यांच्या
आरंभीच्या महत्त्वाच्या निष्कर्षांपैकी एक असा की प्राप्यता
संबंध~$R$ चे अनेक गुणधर्म, उदाहरणार्थ संक्रमणशीलता किंवा
सममिती, योग्य ``मोडल रूपबंधां''च्या साहाय्याने
\emph{मोडल भाषेतच} व्यक्त करता येतात. उदाहरणार्थ,
$R$ ची परावर्तिता ``जर अनिवार्यतः~$!A$, तर~$!A$'' या
रूपबंधाशी ``अनुरूप'' आहे, असे मोडल तर्कशास्त्रज्ञ म्हणतात.
\Ax{D}, \Ax{T}, \Ax{B}, \Ax{4} आणि~\Ax{5} हे रूपबंध
एकत्र करून मिळणाऱ्या अभिजात मोडल तर्कशास्त्रांपैकी
काहींच्या, उदाहरणार्थ \Log{S4} आणि \Log{S5}, अनुरूपता
सिद्धांताचा आपण मुख्यतः अभ्यास करू.

\end{document}
